% 从旧数学附录迁出的高次剩余说明；不新增算法。
\section{素数模高次剩余与单位根}\label{knowledge-higher-residue}
实现见第~\pageref{compact-KthResidue}~页；依赖离散对数、线性同余和快速幂。
设 $p$ 为素数，$g$ 是原根，$k>0$。$a=0$ 时根只有零；否则 $a=g^b$、$x=g^y$，
\[
ky\equiv b\pmod{p-1}.
\]
令 $d=\gcd(k,p-1)$，$d\nmid b$ 时无解；否则 $y=y_0+i(p-1)/d$，$0\le i<d$，存为
\[
\texttt{first}=g^{y_0},\qquad
\texttt{ratio}=g^{(p-1)/d},\qquad \texttt{count}=d.
\]
依次乘 ratio 枚举，不保证递增。$p\le10^{12}$，调用者提供已验证的素数与原根；
$k$ 为 unsigned long long，$a$ 先正规化。空 optional 表示无解，\texttt{\{0,1,1\}} 表示唯一零根。
期望时间、空间 $O(\sqrt p)$，枚举另需 $O(d)$；模乘须扩展到 128 位。

\section{素数幂模高次剩余与根的提升}\label{knowledge-prime-power-residue}
令 $M=p^e$。$a=0$ 时根为 $p^{\lceil e/k\rceil}$ 的倍数。
否则 $t=v_p(a)<e$，须 $k\mid t$，否则无解；取 $s=t/k$，$q=p^{e-t}$，$x=p^s y$，解
\[
y^k\equiv a/p^t\pmod q,
\]
每个单位根 $y$ 对应 $x=p^s y+\ell p^s q$，$0\le\ell<p^{t-s}$。
奇素数用模 $p^e$ 的原根转线性同余；$q=2$ 时单位根只有 1；
$q\ge4$ 为二的幂时单位写成 $(-1)^z5^y$，$z$ 模 2、$y$ 模 $q/4$（$q=4$ 只有符号分支），分别解指数同余。
二的高次幂没有原根。时间由 BSGS 决定，期望 $O(\sqrt q)$。

\section{一般合数模高次剩余：CRT 组合}\label{knowledge-composite-residue}
设 $M=\prod_t m_t$，$m_t=p_t^{e_t}$ 两两互素。分别求局部根族，任一局部无解则整体无解；
整体根数为 $\prod_t c_t$。预处理
\[
E_t=\frac M{m_t}\left(\left(\frac M{m_t}\right)^{-1}\bmod m_t\right)\bmod M,
\qquad x=\sum_t r_tE_t\bmod M.
\]
按混合进制索引可逐个取根，不需建出全部根；CRT 乘积扩展到 128 位。$M=1$ 时唯一代表为 0。

\section{自动准备模数分解与原根}\label{knowledge-residue-setup}
奇素数幂 $m=p^e$ 有原根，设 $\varphi=m/p\cdot(p-1)$，候选 $g$ 为原根当且仅当
\[
p\nmid g,\qquad g^{\varphi/q}\not\equiv1\pmod m
\quad\text{对每个不同素因子 }q\mid\varphi.
\]
须检查模 $p^e$ 的阶，不能只取模 $p$ 的原根。准备成本含 $M$ 与各 $\varphi$ 的分解及逐个候选的模幂；
同一模数多个方程应复用分解结果。
